% \iffalse meta-comment
%
% hit-bibstyle.dtx 是各个类共用的参考文献选项与样式
%
% \fi
%
% \section{共享参考文献样式}
%
% \changes{v3.2a}{2026/08/16}{参考文献分 \texttt{bibtex} 与 \texttt{biber} 两条路线，
%   \option{bib}/\option{backend} 选，各开一个子组（\option{bibtex}/\option{style} 挑
%   \file{.bst}，\option{biber} 原封透传 \pkg{biblatex} 的装载选项）；biblatex 一线由类
%   自己装载配置，不 fork 样式文件，运行时补丁把方括号与逗号冒号分号接到
%   \option{punct} 上；与 natbib 一族的冲突守卫照 thuthesis 用内核钩子写；
%   \option{struct} 的 \option{bibliography} 换算成 \cs{addbibresource}；键挂到
%   \cs{hitsetup} 末尾统一下发}
% \changes{v3.2a}{2026/09/04}{两条线同名同义的键：\option{punct}
%   （\texttt{half}｜\texttt{full}｜\texttt{by-lang}，默认 \texttt{full} 照学校规范的
%   示例）、\option{sentence-case}、\option{show-sl-sn}（默认补「不详」）、
%   \option{show-note}、六个显示开关、\option{comma-in-cite} 与
%   \option{brackets-in-cite}（\cs{bibpunct} 收转接宏，键随设随生效）、
%   \option{correct-punct}（默认关）、\option{uppercase-family}（西文姓氏默认全大写，
%   跟官方底座）、\option{numbering-align}、\option{max-bib-names}}
% \changes{v3.2a}{2026/08/16}{传统一线的键走官方 \pkg{gbt7714} 的控制通道下发，下发前
%   先探键在不在（TeX Live 2025 带的 v2.1.8 没有这条通道），挂到 \cs{bibliography}
%   自己身上}
% \changes{v3.2a}{2026/09/13}{著录列表：两个类相同的那几段抽出来共用，标号与正文的
%   间距抽成 \cs{hit@bib@labelsep:}；条目之间不再留空（范例的著录步进一律 19.5）；
%   人文社科类退回 1.25 倍多行距（正文是固定值 23 磅，照两份人文范例）；三参数版
%   \cs{hit@bib@list:nnn} 给成果页改悬挂与段距}
%
% 两档底下的文献宏包不是一个：终稿走 \pkg{gbt7714}，报告走 \pkg{natbib}。
% 但三个选项、引用样式、\cs{inlinecite}、著录列表的排版参数这几件事两边写的是
% 同一份代码，原先各抄一遍。
%
% 本文件排在共享层，也就是类骨架之前。\pkg{docstrip} 逼的：它给同一个输出里
% 重复出现的输入文件改名（原名后面添空格），并要求每个输出的输入文件表是全局表的
% 子序列。一个输出一旦用上自己独有的文件，扫描位置就落到全局表尾部，再回头引用
% 排在前面的共享文件就找不到了。所以共享文件只能全部排在类骨架第一次出现之前，
% 这也是页面模块必须按类分文件的原因。
%
% 排得早就带来两条限制，都体现在留在各个类里的那几行上：
% \cs{citestyle} 要等文献宏包加载完才能调用，所以那一行留在各个类的文献模块里；
% \cs{bibauthor} 与 \cs{g\_bib\_author\_cs:n} 用 \cs{exp\_args:NNe} 在定义时就把
% \option{bibauthorupper} 的取值展开定死，得排在 \cs{ProcessKeyOptions} 之后，
% 比这里晚，所以两档各留一份。
%
% \subsection{选项}
%
% \DescribeMacro{bibauthorupper}
% \DescribeMacro{bibmaxauthor}
% \DescribeMacro{biblabelalign}
% \option{bibauthorupper} 取 \texttt{full} 时作者姓名全大写；\option{bibmaxauthor}
% 是列表里最多显示几位作者，取 \texttt{0} 表示不限；\option{biblabelalign} 取
% \texttt{left} 时标号左对齐。
%    \begin{macrocode}
%<*shared-bibstyle>

%    \end{macrocode}
%    \begin{macrocode}
\str_new:N \g__hit_bib_cite_punct_str
\str_gset:Nn \g__hit_bib_cite_punct_str { half }
\str_new:N \g__hit_bib_cite_brackets_str
\str_gset:Nn \g__hit_bib_cite_brackets_str { half }
\bool_new:N \g_bib_authorupper_bool
\bool_gset_true:N \g_bib_authorupper_bool
\keys_define:nn { hit / class }
  {
    bibauthorupper .code:n =
      {
        \str_if_eq:nnT {#1} {full}
          { \bool_gset_true:N \g_bib_authorupper_bool }
      } ,
  }

\int_new:N \g_bib_maxauthor_int
\int_gset:Nn \g_bib_maxauthor_int {3}
\int_new:N \g__hit_bib_maxcite_int
\int_gset:Nn \g__hit_bib_maxcite_int { -2 }
\keys_define:nn { hit / class }
  {
    bibmaxauthor .code:n =
      {
        \int_gset:Nn \g_bib_maxauthor_int {#1}
        \int_if_zero:nT { \g_bib_maxauthor_int }
          { \int_gdecr:N \g_bib_maxauthor_int }
      } ,
  }

\str_new:N \g_bib_labelalign_str
\keys_define:nn { hit / class }
  { biblabelalign .code:n = { \str_gset:Nn \g_bib_labelalign_str {#1} } }

%    \end{macrocode}
%
% \subsection{引用样式}
%
% \begin{macro}{\bibstyle@numerical}
% \changes{v3.0.15}{2021/05/06}{删除 inline 等定义，此命令与 \pkg{gbt7714} 冲突}
% 数字上标式引用。\cs{bibstyle@inline} 两档不一样：报告自己定义一份，
% 学位论文用 \pkg{gbt7714} 带的那份，所以它不在这里。
%
% 括号与分隔都从转接宏取：\cs{bibpunct} 存进 \pkg{natbib} 的是宏名，
% 每次引用展开的都是键的现值。
%    \begin{macrocode}
\cs_new:Npn \hit@bib@cite@open:
  { \str_if_eq:VnTF \g__hit_bib_cite_brackets_str { full } {［} {[} }
\cs_new:Npn \hit@bib@cite@close:
  { \str_if_eq:VnTF \g__hit_bib_cite_brackets_str { full } {］} {]} }
\cs_new:Npn \hit@bib@cite@sep:
  {
    \str_case:VnF \g__hit_bib_cite_punct_str
      { { half-with-space } { ,~ } { full } { ，} } { , }
  }
\cs_set_protected:Npn \bibstyle@numerical
  {
    \bibpunct { \hit@bib@cite@open: } { \hit@bib@cite@close: }
      { \hit@bib@cite@sep: } {s} {,} { \textsuperscript {,} }
  }
%    \end{macrocode}
% \end{macro}
%
% \begin{macro}{\inlinecite,\onlinecite}
% \changes{v2.0.06}{2018/12/05}{在 \cs{inlinecite} 内添加空格}
% \cs{inlinecite}\marg{键} 把引用排成正文里的样子，不上标。
%    \begin{macrocode}
\DeclareRobustCommand\inlinecite{\@inlinecite}
\cs_set:Npn \@inlinecite #1
  {
    \group_begin:
      \citestyle { inline }
      \cs_set_eq:NN \@cite \NAT@citenum
      \citep {#1}
    \group_end:
  }
\cs_set_eq:NN \onlinecite \inlinecite
%    \end{macrocode}
% \end{macro}
%
% \subsection{著录列表}
%
% \begin{macro}{\hit@bib@list:nn}
% \cs{hit@bib@list:nn}\marg{最宽的标号}\marg{标号与正文的间距} 开一个著录列表。
% 两档的 \env{thebibliography} 前后差得多，标题怎么排、分页惩罚怎么给都不一样，
% 中间这段 \cs{list} 的参数是一样的，只有间距要各传各的。
%
% 条目之间不留空。三份范例的文献表逐行量，条目内、条目间的基线步进一律
% 19.5（1.25 倍多行距的行盒），先前 \cs{itemsep} 从字号默认里减 0.8\,em
% 剩下的那 0.35 与 \cs{parsep} 都是多出来的，两条后端线一并清零。
%
% 三参数版多收一段列表参数，排在公共参数之后执行，调用方要改哪一个就覆盖
% 哪一个。成果页的列表与参考文献表机制同源、有两个数不同（悬挂宽与段前
% 段后，范例如此，见 \file{hit-final-backmatter.dtx}），就从这儿进。
%
% \emph{那个间距抽出来。}它先前写死在四个调用点上（两档各有传统与
% \pkg{biblatex} 两条线），终稿还带着一句“深圳校区取 0”的判断，抄了两遍。
% 成果页也要按同一套排编号，再抄一遍就是第五处。改成一个取值宏
% \cs{hit@bib@labelsep:}，值存在 \cs{g\_\_hit\_bib\_labelsep\_tl} 里，
% 默认半个字宽；标号排全角方括号的那几档由终稿在选项处理完之后改写
% 成零（范例如此，见 \file{hit-final-bib.dtx}），报告不分校区所以照旧。
%
% \emph{存的是式子不是长度。}\texttt{0.5em} 里的 em 跟着当时的字体走，
% 存进长度寄存器就等于在类加载那一刻定死（那时还是十点的默认字号，
% 折出来 5\,pt，而著录列表里是小四的 6\,pt）。所以存在 token 列表里，
% 到 \cs{labelsep} 那一句才求值。
%    \begin{macrocode}
\tl_new:N \g__hit_bib_labelsep_tl
\tl_gset:Nn \g__hit_bib_labelsep_tl { 0.5em }
\cs_new:Npn \hit@bib@labelsep: { \g__hit_bib_labelsep_tl }
%    \end{macrocode}
%
% \emph{列表参数之外还有几项。}断行放松、句点后不加英文的句末空格、
% 分页惩罚——这三样先前写在 \env{thebibliography} 的正文里，成果页那一族
% 抄不到，于是同一批著录在两处的字距不一样：句点接西文时成果页多出一段
% 句末空格，跨 \TeX~Live 版本还各多一点（实测「\texttt{.S}」两版差
% 2.37\,bp）。抽成一句，两处共用。
%
% \option{split-entry} 决定条目能不能跨页断开：开着就把孤行寡行的惩罚放到
% 最松，关着则一律 4000 顶住。
%    \begin{macrocode}
\cs_new_protected:Npn \hit@bib@list@rules:
  {
    \sloppy \frenchspacing
%    \end{macrocode}
%
% 人文社科类的正文是固定值 23 磅（见 \file{hit-defaults.dtx}），文献表与成果
% 条目却不是：两份人文范例里这两处的段落都是 \texttt{line=300 auto}，Word
% 导出逐行量步进 19.4～19.5，与理工一个数。所以列表里把撑杆的固定值档关掉、
% \cs{baselineskip} 退回小四宋体的 1.25 倍行盒（12.18 加 7.28）。标志位只在
% 终稿里有。
%    \begin{macrocode}
    \bool_lazy_and:nnT
      { \bool_if_exist_p:N \g__hit_hass_bool } { \bool_if_p:N \g__hit_hass_bool }
      {
        \bool_set_false:N \l_docxlike_format_strut_exact_bool
        \dim_set:Nn \baselineskip { 19.46bp }
      }
    \bool_if:NTF \g__hit_splitbibitem_bool
      {
        \int_set:Nn \clubpenalty { 0 }
        \int_set_eq:NN \@clubpenalty \clubpenalty
        \int_set:Nn \widowpenalty { 0 }
        \int_set:Nn \interlinepenalty { -50 }
      }
      {
        \int_set:Nn \clubpenalty { 4000 }
        \int_set_eq:NN \@clubpenalty \clubpenalty
        \int_set:Nn \widowpenalty { 4000 }
        \int_set:Nn \interlinepenalty { 4000 }
      }
    \sfcode `\. = \@m \scan_stop:
  }
%    \end{macrocode}
%
%    \begin{macrocode}
\cs_new_protected:Npn \hit@bib@list:nnn #1#2#3
  {
    \list { \@biblabel { \@arabic \c@enumiv } }
      {
        \str_if_eq:VnT \g_bib_labelalign_str { left }
          { \RenewDocumentCommand \makelabel { m } { ##1 \hfill } }
        \settowidth { \labelwidth } { \@biblabel {#1} }
        \dim_set:Nn \labelsep {#2}
        \dim_set:Nn \itemindent { 0pt }
        \dim_set:Nn \leftmargin { \labelsep + \labelwidth }
        \skip_zero:N \itemsep
        \skip_zero:N \parsep
        \usecounter { enumiv }
        \cs_set_eq:NN \p@enumiv \@empty
        \cs_set:Npn \theenumiv { \@arabic \c@enumiv }
        #3
      }
  }
\cs_new_protected:Npn \hit@bib@list:nn #1#2 { \hit@bib@list:nnn {#1} {#2} { } }
%    \end{macrocode}
% \end{macro}
%
% \subsection{后端}
%
% \DescribeMacro{backend}
% \option{bib}/\option{backend} 选参考文献走哪条路线，取值按「要跑哪个程序」起名：
% \texttt{bibtex} 是传统一线（\pkg{gbt7714}/\pkg{natbib} 加 \file{.bst}，默认），
% \texttt{biber} 是 \pkg{biblatex} 一线（样式走 \pkg{biblatex-gb7714-2015}）；
% \texttt{biblatex} 当 \texttt{biber} 的别名也认。默认 \texttt{auto}——没写就走
% 传统一线，但用户要是自己在导言区装了 \pkg{biblatex}，模板会侦测到并且不再装
% \pkg{gbt7714}/\pkg{natbib}（侦测在导言区末尾，见各类的 bib 模块），效果等同
% \texttt{biber}。这一招照抄 thuthesis。
%
% \pkg{biblatex} 的装载不能拖到导言区末尾（它自己就在装载时挂
% \texttt{begindocument} 一族的钩子，晚了赶不上趟），所以在 \option{backend}
% 所在的那组 \option{bib}\texttt{=\{...\}} 处理完的当口就装。这样
% \option{style}、\option{author-case} 这些「装载时就得定死」的键只要跟
% \option{backend} 写在同一组里，先后随意；分开写则必须先键后 \option{backend}，
% 写反了有中文提示（见 \file{hit-messages.dtx} 的 bib-key-locked 一族）。
%    \begin{macrocode}
%    \end{macrocode}
%
% \begin{macro}{punct}
% \option{punct} 管著录表的标点宽度，三档：\texttt{full}（默认，不分语种
% 一律全角逗号冒号圆括号方括号，句点恒半角带空格，学校规范的示例与旧深圳版
% 都是这个排法）、\texttt{half}（一律半角，旧哈尔滨版的排法）、
% \texttt{by-lang}（按条目语言，中文全角西文半角）。传统一线映射官方
% \file{.bst} 的 \texttt{bibpunct} 三档（\texttt{GB}/\texttt{half}/%
% \texttt{bylanguage}，控制通道在 \cs{bibliography} 当口才读）；biblatex
% 一线的方括号由下面的补丁按此键分发。两条线都随设随生效，不锁时机。
%    \begin{macrocode}
\str_new:N \g__hit_bib_punct_str
\str_gset:Nn \g__hit_bib_punct_str { full }
\keys_define:nn { hit / bib }
  {
    punct .choices:nn = { half , full , by-lang }
      { \str_gset:NV \g__hit_bib_punct_str \l_keys_choice_tl } ,
    punct .value_required:n = true ,
  }
%    \end{macrocode}
% \end{macro}
%
% \begin{macro}{sentence-case}
% \option{sentence-case} 管西文题名要不要转成句首大写（只首字母大写，其余小写），
% 不转就照 \file{.bib} 里写的原样出。三个字段各一档：
% \begin{latex}
% bib = { sentence-case = { title = true , booktitle = true , journal = false } }
% \end{latex}
% 也可以一次设三个：\option{sentence-case}\texttt{=true|false|auto}。
% \texttt{auto} 退回规范默认，就是上面那一行的取值——题名与析出文献题名转，
% 期刊名不转（期刊名的缩写形式本来就带自己的大小写，转了会毁掉）。
%
% 名字取自两条后端线共有的说法：\file{.bst} 那边的控制字段叫
% \texttt{CTL\_sentence\_case\_title}，\pkg{biblatex} 那边的命令叫
% \cs{MakeSentenceCase}。不叫 \texttt{title-case}——那在西文排版里指
% 「每个实词首字母大写」，跟这里要的正好相反。
%
% 两线各自落实。\file{.bst} 一线走控制通道下发三个字段，见
% \cs{hit@bib@forward@ctl:}；biblatex 一线改写
% \pkg{biblatex-gb7714-2015} 留的 \cs{mktitlecase} 三个钩子，钩子里读标志位，
% 所以也是随设随生效。
%    \begin{macrocode}
\bool_new:N \g__hit_bib_sentence_title_bool
\bool_new:N \g__hit_bib_sentence_booktitle_bool
\bool_new:N \g__hit_bib_sentence_journal_bool
\bool_gset_true:N  \g__hit_bib_sentence_title_bool
\bool_gset_true:N  \g__hit_bib_sentence_booktitle_bool
\bool_gset_false:N \g__hit_bib_sentence_journal_bool
%    \end{macrocode}
%
% \meta{字段}\marg{取值}\marg{auto 档的规范默认}：规范默认写在调用点上，
% 不另立一张表，省得两处对不上。
%    \begin{macrocode}
\cs_new_protected:Npn \__hit_bib_sentence_set:nnn #1#2#3
  {
    \str_case:nnF {#1}
      {
        { true }  { \bool_gset_true:c  { g__hit_bib_sentence_#2_bool } }
        { false } { \bool_gset_false:c { g__hit_bib_sentence_#2_bool } }
        { auto }
          {
            \str_if_eq:nnTF {#3} { true }
              { \bool_gset_true:c  { g__hit_bib_sentence_#2_bool } }
              { \bool_gset_false:c { g__hit_bib_sentence_#2_bool } }
          }
      }
      { }
  }
\cs_new_protected:Npn \__hit_bib_sentence_all:n #1
  {
    \__hit_bib_sentence_set:nnn {#1} { title }     { true }
    \__hit_bib_sentence_set:nnn {#1} { booktitle } { true }
    \__hit_bib_sentence_set:nnn {#1} { journal }   { false }
  }
\keys_define:nn { hit / bib }
  {
    sentence-case .code:n =
      {
        \str_case:nnF {#1}
          {
            { true }  { \__hit_bib_sentence_all:n { true } }
            { false } { \__hit_bib_sentence_all:n { false } }
            { auto }  { \__hit_bib_sentence_all:n { auto } }
          }
          { \hit@keys@set:nn { hit / bib / sentence-case } {#1} }
      } ,
    sentence-case .default:n = true ,
  }
\keys_define:nn { hit / bib / sentence-case }
  {
    title .choices:nn = { true , false , auto }
      {
        \exp_args:NV \__hit_bib_sentence_set:nnn
          \l_keys_choice_tl { title } { true }
      } ,
    title .default:n = true ,
    booktitle .choices:nn = { true , false , auto }
      {
        \exp_args:NV \__hit_bib_sentence_set:nnn
          \l_keys_choice_tl { booktitle } { true }
      } ,
    booktitle .default:n = true ,
    journal .choices:nn = { true , false , auto }
      {
        \exp_args:NV \__hit_bib_sentence_set:nnn
          \l_keys_choice_tl { journal } { false }
      } ,
    journal .default:n = true ,
  }
%    \end{macrocode}
% \end{macro}
%
%    \begin{macrocode}
\str_new:N \g__hit_bib_backend_str
\str_gset:Nn \g__hit_bib_backend_str { auto }
\tl_new:N \g__hit_bib_style_used_tl
\keys_define:nn { hit / bib }
  {
    backend .choices:nn = { auto , bibtex , biber , biblatex }
      {
        \str_gset:NV \g__hit_bib_backend_str \l_keys_choice_tl
        \str_if_eq:VnT \g__hit_bib_backend_str { biblatex }
          { \str_gset:Nn \g__hit_bib_backend_str { biber } }
      } ,
    backend .value_required:n = true ,
  }
%    \end{macrocode}
%
% \begin{macro}{\hit@bib@backend@apply:}
% 每处理完一次 \cs{hitsetup} 就照 \option{backend} 的现值行动一次。
%
% \emph{挂在 \cs{hitsetup} 末尾，不能挂在根族的 \option{bib} 路由上。} l3keys 认
% \option{bib}/\option{backend} 这种写路径的形式，它直接送进 \texttt{hit\,/\,bib}
% 族，根族那一层的 \option{bib} 键根本不经过，挂在那里的动作也就不执行。用户
% 在导言区写 \cs{hitsetup}\{bib/backend=biber\} 会照编不误，然后收到一句
% 「设晚了」——因为装载动作被推到了别处触发，那时导言期已经过了。
%
% 也不能挪进 \option{backend} 那个键自己的键码：整组 \option{bib}\texttt{=\{...\}}
% 里 \option{biber}\texttt{=\{...\}} 那些子选项要等读完才齐，在 \option{backend}
% 当场装载会漏掉它们。挂在 \cs{hitsetup} 末尾两头都顾上：整组读完了，
% 而且不管用哪种写法进来都会经过。重复调用幂等。
%
% 行动是：\texttt{biber} 而 \pkg{biblatex} 未装就装；
% \texttt{bibtex} 而 \pkg{biblatex} 已装是矛盾，报错；\texttt{auto} 不动。
% 重复设同一个值幂等。
%    \begin{macrocode}
\cs_new_protected:Npn \hit@bib@backend@apply:
  {
    \str_case:Vn \g__hit_bib_backend_str
      {
        { biber }
          { \@ifpackageloaded { biblatex } { } { \hit@bib@load@biblatex: } }
        { bibtex }
          {
            \@ifpackageloaded { biblatex }
              { \msg_error:nn { hithesis } { bib-backend-conflict } } { }
          }
        { auto } { }
      }
  }
%    \end{macrocode}
% \end{macro}
%
% \begin{macro}{\hit@bib@load@biblatex:}
% 装载 \pkg{biblatex}。样式锁死 \pkg{biblatex-gb7714-2015} 的
% \texttt{gb7714-2015}：规范就是国标，biblatex 一线不开挑样式的口子——
% \file{.bst} 一线让挑（\option{bibtex}/\option{style}）是因为 \hologo{BibTeX}
% 能力有限，得靠换 \file{.bst} 微调；biblatex 自己的选项够用，细调走
% \option{biber}\texttt{=\{...\}} 透传。
%
% 有一道时机护栏：demo 的主文件把 \cs{hitsetup} 全写在
% \cs{begin}\marg{document} 之后，元信息那样设没问题，可 \pkg{biblatex}
% 过了导言区就装不进去了。设晚了给一条说清怎么挪的警告然后照旧走传统一线，
% 不让内核那句 Can~be~used~only~in~preamble 把用户晾在原地。
%    \begin{macrocode}
\bool_new:N \g__hit_bib_preamble_bool
\bool_gset_true:N \g__hit_bib_preamble_bool
\AddToHook { begindocument }
  { \bool_gset_false:N \g__hit_bib_preamble_bool }
\cs_new_protected:Npn \hit@bib@load@biblatex:
  {
    \bool_if:NTF \g__hit_bib_preamble_bool
      {
        \use:e
          {
            \exp_not:N \RequirePackage
              [
                backend = biber , style = gb7714-2015
                \tl_if_empty:NF \g__hit_bib_biber_opts_tl
                  { , \exp_not:V \g__hit_bib_biber_opts_tl }
              ] { biblatex }
          }
      }
      { \msg_warning:nn { hithesis } { bib-backend-in-body } }
  }
%    \end{macrocode}
% \end{macro}
%
% \begin{macro}{\hit@bib@key@guard:n}
% 「设晚了」的守卫。\option{uppercase-family} 这些键在 \pkg{biblatex} 装载时换算成
% 它的选项，装载之后再设不生效，键码前面搭一句守卫提示。
% 传统一线不受影响——那边的取值到排版时才读。
%    \begin{macrocode}
\cs_new_protected:Npn \hit@bib@key@guard:n #1
  {
    \@ifpackageloaded { biblatex }
      { \msg_warning:nnn { hithesis } { bib-key-locked } {#1} } { }
  }
%    \end{macrocode}
% \end{macro}
%
% \begin{macro}{bibtex,biber}
% 后端专属的配置各归各的子组。\option{bibtex}/\option{style} 是
% \cs{bibliographystyle} 的参数（\texttt{hithesis}（默认）或
% \hologo{TeX} 树里任何 \file{.bst} 名），排版时才读，什么时候设都行。
% \option{biber} 收一张键值表，原封并进 \pkg{biblatex} 的装载选项——
% \pkg{biblatex-gb7714-2015} 的几十个选项（\texttt{gbpub}、\texttt{gbtype}、
% \texttt{gbannote}……）都从这儿过，模板不逐个转译；装载选项自然要在装载前
% 给，跟 \option{backend} 同组或在其之前，设晚了同一句守卫提示。
%    \begin{macrocode}
\tl_new:N \g__hit_bib_biber_opts_tl
\keys_define:nn { hit / bib }
  {
    bibtex .code:n = { \hit@keys@set:nn { hit / bib / bibtex } {#1} } ,
    biber .code:n =
      {
        \hit@bib@key@guard:n { biber }
        \tl_gset:Nn \g__hit_bib_biber_opts_tl {#1}
      } ,
  }
\keys_define:nn { hit / bib / bibtex }
  { style .code:n = { \tl_gset:Nn \g__hit_bib_style_tl {#1} } }
%    \end{macrocode}
% \end{macro}
%
% \begin{macro}{show-sl-sn}
% \option{show-sl-sn}（\texttt{true}／\texttt{false}）：文献缺出版地或出版者时补不补
% ［出版地不详］（sine~loco）与［出版者不详］（sine~nomine）。国标要补，
% 默认 \texttt{true}。名字取自 kebab-biblatex-gb7714 的同义键。
%
% 落实分两线：biblatex 一线在装载时换算成 \pkg{biblatex-gb7714-2015} 的
% \texttt{gbpub}——注意那边默认是 \texttt{false}（不补），跟 \file{.bst}
% 一线的「按国标补」对不上，所以这里无论用户设没设都显式发一次，两线默认
% 才一致。
%
% 传统一线走官方 gbt7714~v3 的 \texttt{unknownpublisher}（控制通道，见
% \cs{hit@bib@forward@ctl:}），两线都真实可开可关。
%    \begin{macrocode}
\bool_new:N \g__hit_bib_slsn_bool
\bool_gset_true:N \g__hit_bib_slsn_bool
\keys_define:nn { hit / bib }
  {
    show-sl-sn .choices:nn = { true , false }
      {
        \hit@bib@key@guard:n { show-sl-sn }
        \str_if_eq:VnTF \l_keys_choice_tl { true }
          { \bool_gset_true:N  \g__hit_bib_slsn_bool }
          { \bool_gset_false:N \g__hit_bib_slsn_bool }
      } ,
    show-sl-sn .default:n = true ,
  }
%    \end{macrocode}
% \end{macro}
%
% \begin{macro}{show-mark,show-medium,link-title,show-url,show-doi,show-eprint}
% 六个布尔开关，两边的官方机制一一对应（\file{.bst} 的 CTL 键 ↔
% \pkg{biblatex} 的 gb 系或核心选项）：\option{show-mark}（类型标 ［J］，
% \texttt{entrytypeid}↔\texttt{gbtype}）、\option{show-medium}（载体标
% \texttt{/OL}，\texttt{entrymediumid}↔\texttt{gbmedium}）、
% \option{link-title}（题名加超链接，\texttt{linktitle}↔\texttt{gbtitlelink}）、
% \option{show-url}/\option{show-doi}/\option{show-eprint}
% （\texttt{url}/\texttt{doi}/\texttt{eprint}，biblatex 核心同名）。
% 默认取 \file{.bst} 的内置默认：前两个开、后面 \texttt{link-title} 关、
% \texttt{url}/\texttt{doi} 开、\texttt{eprint} 关——biblatex 核心的
% \texttt{eprint} 本来默认开，跟 \file{.bst} 扯平成关。
%    \begin{macrocode}
\bool_new:N \g__hit_bib_showmark_bool
\bool_gset_true:N \g__hit_bib_showmark_bool
\bool_new:N \g__hit_bib_showmedium_bool
\bool_gset_true:N \g__hit_bib_showmedium_bool
\bool_new:N \g__hit_bib_linktitle_bool
\bool_new:N \g__hit_bib_showurl_bool
\bool_gset_true:N \g__hit_bib_showurl_bool
\bool_new:N \g__hit_bib_showdoi_bool
\bool_gset_true:N \g__hit_bib_showdoi_bool
\bool_new:N \g__hit_bib_showeprint_bool
\bool_new:N \g__hit_bib_shownote_bool
\cs_new_protected:Npn \__hit_bib_show_key:nN #1#2
  {
    \keys_define:nn { hit / bib }
      {
        #1 .choices:nn = { true , false }
          {
            \hit@bib@key@guard:n {#1}
            \str_if_eq:VnTF \l_keys_choice_tl { true }
              { \bool_gset_true:N #2 }
              { \bool_gset_false:N #2 }
          } ,
        #1 .default:n = true ,
      }
  }
\__hit_bib_show_key:nN { show-mark }   \g__hit_bib_showmark_bool
\__hit_bib_show_key:nN { show-medium } \g__hit_bib_showmedium_bool
\__hit_bib_show_key:nN { link-title }  \g__hit_bib_linktitle_bool
\__hit_bib_show_key:nN { show-url }    \g__hit_bib_showurl_bool
\__hit_bib_show_key:nN { show-doi }    \g__hit_bib_showdoi_bool
\__hit_bib_show_key:nN { show-eprint } \g__hit_bib_showeprint_bool
%    \end{macrocode}
%
% \option{show-note} 管 \file{.bib} 里的 \texttt{note} 字段排不排在条目末尾。
% 默认关，与 \file{.bst} 的出厂档一致：参考文献表照国标著录，不添自己的话。
% 成果页那边常要在条目后缀一句“SCI 收录，IDS 号为……”“EI 收录号：……”，
% 那种就写进 \texttt{note} 再把这个键打开。\file{.bst} 侧的开关是
% \texttt{CTL\_note}，本来就有，先前没接出来。
%    \begin{macrocode}
\__hit_bib_show_key:nN { show-note }   \g__hit_bib_shownote_bool
%    \end{macrocode}
% \end{macro}
%
% \begin{macro}{cite-punct,correct-punct}
% \option{comma-in-cite} 管正文引用序号之间的分隔：\texttt{half}（默认，
% \texttt{[3,5]}）、\texttt{half-with-space}（\texttt{[3,~5]}）、
% \texttt{full}（全角逗号）。\option{brackets-in-cite} 管引用的方括号：
% \texttt{half}（默认）或 \texttt{full}（［3,5］）。传统一线经
% \cs{bibstyle@numerical} 里的转接宏落到 \pkg{natbib} 上；biblatex 一线
% 分隔落在 gb 系 cbx 的三个分隔宏、方括号落在 \cs{gbleftbracket} 的引用侧
% 分支（见 \cs{hit@bib@patch@punct:}），全部排版时读值。两条线都不锁时机。
%
% \option{correct-punct}（默认关）管数据清洗：字段里写的标点按
% \option{punct} 的档位归一（全角档把 \texttt{,~:~;~!~?~()} 换成全角，
% 半角档反向）。传统一线是 \file{.bst} 的 \texttt{convertpunct}；biblatex
% 一线用 biber 的 sourcemap 镜像同一套替换（见
% \cs{hit@bib@declare@sourcemaps:}），姓名字段除外——author 里的半角逗号是
% 姓名分隔语法，装载前动它会毁解析，这一点与 \file{.bst}（排版期才转）不同。
%    \begin{macrocode}
\bool_new:N \g__hit_bib_correct_punct_bool
\bool_gset_false:N \g__hit_bib_correct_punct_bool
\keys_define:nn { hit / bib }
  {
    comma-in-cite .choices:nn = { half , half-with-space , full }
      { \str_gset:NV \g__hit_bib_cite_punct_str \l_keys_choice_tl } ,
    comma-in-cite .value_required:n = true ,
    brackets-in-cite .choices:nn = { half , full }
      { \str_gset:NV \g__hit_bib_cite_brackets_str \l_keys_choice_tl } ,
    brackets-in-cite .value_required:n = true ,
    correct-punct .choices:nn = { true , false }
      {
        \hit@bib@key@guard:n { correct-punct }
        \str_if_eq:VnTF \l_keys_choice_tl { true }
          { \bool_gset_true:N  \g__hit_bib_correct_punct_bool }
          { \bool_gset_false:N \g__hit_bib_correct_punct_bool }
      } ,
    correct-punct .default:n = true ,
  }
%    \end{macrocode}
% \end{macro}
%
% \begin{macro}{\hit@bib@print:,\hit@bib@forward@ctl:}
% 排文献表的总出口，两档的 \cs{hitbackmatter} 流程都走它：\pkg{biblatex}
% 装了就 \cs{printbibliography}，没装就是原来那两句 \cs{bibliographystyle} 加
% \cs{bibliography}。判据直接看宏包在不在，不看 \option{backend} 的字面值，
% 免得「用户自己装了 biblatex 而没写键」这一档跟状态记录对不上。
%
% \emph{下发前先探键在不在。} 这条通道是官方 \pkg{gbt7714} v3（2026-07）才有的。
% TeX Live 2025 带的是 v2.1.8，\cs{@ifpackageloaded} 照样为真，但族里没有这些键，
% 12 个键会各报一次「Can be used only in preamble」，键名还被原样排进文档
% （v2 那边同名的是 \pkg{natbib} 的 \cs{bibpunct} 一族，都是导言区专用命令，
% 禁用之后参数没人吃，落进正文）。所以改成探 \option{bibpunct} 在不在：
% 在就下发，不在就整段跳过，键不生效但文档照排。
%
% 传统一线排之前先把键下发给官方通道：\pkg{gbt7714.sty} 把这些选项写成一个
% \texttt{-BSTcontrol} 控制文件、挂进 \cs{bibliography} 的库表，官方底座的
% \file{.bst} 从那里读。控制文件在 \cs{bibliography} 执行的当口才写，所以这些
% 键在传统一线什么时候设都来得及，不用锁。\texttt{0}（不限）换成 999。
% 报告只装 \pkg{natbib} 没有这条通道，吃 \file{.bst} 内置默认。
%    \begin{macrocode}
\tl_const:Nn \c__hit_bib_ctl_suffix_tl { -BSTcontrol }
\tl_const:Nn \c__hit_bib_ctl_type_tl { GBT7714BSTCTL }
\tl_const:Nn \c__hit_bib_ctl_key_tl { GBT7714:BSTcontrol }
\iow_new:N \l__hit_bib_ctl_iow
\prop_new:N \l__hit_bib_ctl_prop
\cs_new_protected:Npn \__hit_bib_ctl_bool:nN #1#2
  { \prop_put:Nnx \l__hit_bib_ctl_prop {#1} { \bool_if:NTF #2 { true } { false } } }
\cs_new_protected:Npn \__hit_bib_ctl_names:nN #1#2
  {
    \prop_put:Nnx \l__hit_bib_ctl_prop {#1}
      { \int_compare:nNnTF #2 < 0 { 999 } { \int_use:N #2 } }
  }
\cs_new_protected:Npn \hit@bib@forward@ctl:
  {
    \prop_clear:N \l__hit_bib_ctl_prop
    \prop_put:Nnx \l__hit_bib_ctl_prop { CTL_bib_punct }
      {
        \str_case:VnF \g__hit_bib_punct_str
          { { half } { half } { by-lang } { bylanguage } } { GB }
      }
    \__hit_bib_ctl_bool:nN { CTL_convert_punct }     \g__hit_bib_correct_punct_bool
    \__hit_bib_ctl_bool:nN { CTL_entry_type_id }     \g__hit_bib_showmark_bool
    \__hit_bib_ctl_bool:nN { CTL_entry_medium_id }   \g__hit_bib_showmedium_bool
    \__hit_bib_ctl_bool:nN { CTL_link_title }        \g__hit_bib_linktitle_bool
    \__hit_bib_ctl_bool:nN { CTL_url }               \g__hit_bib_showurl_bool
    \__hit_bib_ctl_bool:nN { CTL_doi }               \g__hit_bib_showdoi_bool
    \__hit_bib_ctl_bool:nN { CTL_eprint }            \g__hit_bib_showeprint_bool
    \__hit_bib_ctl_bool:nN { CTL_note }              \g__hit_bib_shownote_bool
    \__hit_bib_ctl_bool:nN { CTL_uppercase_family }  \g_bib_authorupper_bool
    \__hit_bib_ctl_bool:nN { CTL_sentence_case_title }
      \g__hit_bib_sentence_title_bool
    \__hit_bib_ctl_bool:nN { CTL_sentence_case_booktitle }
      \g__hit_bib_sentence_booktitle_bool
    \__hit_bib_ctl_bool:nN { CTL_sentence_case_journal }
      \g__hit_bib_sentence_journal_bool
    \__hit_bib_ctl_bool:nN { CTL_unknown_publisher } \g__hit_bib_slsn_bool
    \__hit_bib_ctl_names:nN { CTL_max_bib_names } \g_bib_maxauthor_int
    \__hit_bib_ctl_names:nN { CTL_min_bib_names } \g_bib_maxauthor_int
    \int_compare:nNnF \g__hit_bib_maxcite_int = { -2 }
      {
        \__hit_bib_ctl_names:nN { CTL_max_cite_names } \g__hit_bib_maxcite_int
        \__hit_bib_ctl_names:nN { CTL_min_cite_names } \g__hit_bib_maxcite_int
      }
    \iow_open:Nn \l__hit_bib_ctl_iow
      { \jobname \c__hit_bib_ctl_suffix_tl .bib }
      \iow_now:Nx \l__hit_bib_ctl_iow
        {
          @ \c__hit_bib_ctl_type_tl \c_left_brace_str
          \c__hit_bib_ctl_key_tl ,
        }
      \prop_map_inline:Nn \l__hit_bib_ctl_prop
        {
          \iow_now:Nx \l__hit_bib_ctl_iow
            { \c_space_tl \c_space_tl ##1 \c_space_tl = \c_space_tl {##2} , }
        }
      \iow_now:Nx \l__hit_bib_ctl_iow { \c_right_brace_str }
    \iow_close:N \l__hit_bib_ctl_iow
    \legacy_if:nT { @filesw }
      {
        \iow_now:Nx \@auxout
          { \token_to_str:N \citation { \c__hit_bib_ctl_key_tl } }
      }
  }
\cs_new_protected:Npn \hit@bib@print:
  {
    \@ifpackageloaded { biblatex }
      { \printbibliography }
      {
        \exp_args:NV \bibliographystyle \g__hit_bib_style_tl
        \exp_args:Nv \bibliography { g__hit_structure_bibliography_tl }
      }
  }
%    \end{macrocode}
%
% \emph{下发要挂在 \cs{bibliography} 上，不能只挂在 \cs{hit@bib@print:} 里。}
% v3.1 时代的文档在正文里直接写 \cs{bibliographystyle}\marg{hithesis} 加
% \cs{bibliography}\marg{reference}，走不到 \cs{hit@bib@print:}。下发只写在那里的话，
% 这类文档编得过、但 \option{bib} 底下的键一个都不落地，用户也收不到任何提示——
% 实测老文档里 \option{max-bib-names}\texttt{=1} 无效，同一个键在新范例上是好的。
%
% 所以改成包住 \cs{bibliography}：谁调都先下发一次。\cs{hit@bib@print:} 里那一句
% 因此撤掉，两条路合成一条。下发本身是 \cs{keys\_set:nn}，重复调用无害。
%
% 挂的时机是 \cs{AtBeginDocument}：\pkg{natbib} 与 \pkg{gbt7714} 都在导言区加载，
% 太早包会包到还没被它们改过的那一版。\pkg{biblatex} 一线没有 \cs{bibliography}
% 这条路，所以只在传统一线包。
%    \begin{macrocode}
\AtBeginDocument
  {
    \@ifpackageloaded { biblatex } { }
      {
        \cs_if_exist:NT \bibliography
          {
            \cs_set_eq:NN \hit@bib@inner@bibliography \bibliography
            \cs_set_protected:Npn \bibliography #1
              {
                \hit@bib@forward@ctl:
                \use:e
                  {
                    \exp_not:N \hit@bib@inner@bibliography
                      { \jobname \c__hit_bib_ctl_suffix_tl , \exp_not:n {#1} }
                  }
              }
          }
      }
  }
%    \end{macrocode}
% \end{macro}
%
% \begin{macro}{\hit@bib@biblatex@common:}
% \pkg{biblatex} 装载完立刻要做的、两档相同的那部分（挂在下面那个
% \texttt{package/biblatex/after} 钩子上，类专属的排版配置由各类的
% \cs{hit@bib@biblatex@setup:} 接着做）。
%
% 三个选项键换算成 \pkg{biblatex} 的选项：\option{author-case} 对
% \texttt{gbnamefmt}（\texttt{full} 是全大写，即国标默认；不设则保持
% \file{.bib} 里写的样子，与 \file{.bst} 一线的默认一致）、
% \option{maximum-author} 对 \texttt{maxbibnames}/\texttt{minbibnames}
% （\texttt{0} 表示不限，换成 999）、\option{label-align} 对
% \texttt{gbalign}。标点宽度 \option{punct} 在 biblatex 一线全靠
% 运行时补丁：方括号圆括号走 \cs{bibleftbracket} 一族的领口，逗号冒号分号
% 走 \cs{abx@puncthook} 单点（2015 样式没有 2025 代的
% \texttt{gbpunctwidth} 旋钮，能力用补丁补齐），见下方两个 patch 宏。
% gb 系选项认准 \cs{mkgbnumlabel} 在不在——用户点名
% 别家样式（比如 apa）时这些选项不存在，塞过去要报错。
%
% 这些换算必须在装载当口做：gb 系样式把选项码挂进了每个条目的选项链，
% 排版时逐条重放，导言区后半再 \cs{ExecuteBibliographyOptions} 会被
% 重放盖掉（实测）。所以对应的键设晚了只能提示，见上面的守卫。
%
% \cs{inlinecite} 的接法：gb7714-2015 的 cbx 自己用 \cs{newcommand} 定义一个
% \cs{inlinecite}（就是 \cs{parencite}），类里先定义的那份会让它装载报错，
% 所以装载前先把类的定义撤掉、让 cbx 的进来；点名的样式没带
% \cs{inlinecite} 的（比如 apa），装完再把 robust 外壳补回去，里面换成
% \cs{parencite}。\cs{onlinecite} 重新对齐到成品上。\cs{supercite} 也归一：
% \pkg{biblatex} 核心那份是裸上标数字不带括号，改指 gb 系 cbx 的 \cs{cite}
% （上标带方括号），与传统一线补的别名同义。
% \cs{bibliographystyle} 降级成一条提示，老文档里残留这一句不至于炸。
%
% 正文里的老写法也要接住：报告的示例正文就是两句
% \cs{bibliographystyle} 加 \cs{bibliography}。\pkg{biblatex} 一线把正文的
% \cs{bibliography} 映射成 \cs{printbibliography}（文献库那时早该注册完了，
% 参数吞掉不用），\cs{bibliographystyle} 无论导言区还是正文都降级成提示。
% 映射挂 \texttt{begindocument/end}：\pkg{biblatex} 把这两个名字登进内核的
% \cs{@onlypreamble} 表，内核在 \cs{document} 收尾处统一禁用
% （\cs{@preamblecmds}，晚于 \texttt{begindocument} 钩子），要排在那之后
% 才盖得回来。
%
% \begin{macro}{\hit@bib@patch@brackets:}
% biblatex 一线的方括号补丁。gb7714-2015 的 bbx 把著录里所有方括号（类型标、
% 访问日期、序号标签的缺省形）都收口在 \pkg{biblatex} 内核的
% \cs{bibleftbracket}/\cs{bibrightbracket} 一对宏上，补丁只动这两个名字，
% 不动样式文件。分发看 \option{punct} 的现值：\texttt{full} 全角、
% \texttt{half} 走存下来的原定义、\texttt{by-lang} 按条目语言（bbx 用
% biber 的 sourcemap 把语言标进 \texttt{userd} 域，排版时逐条可查）。
% 现值是排版时才读的，键随设随生效。引用标签的方括号（\cs{gbleftbracket}
% 一族）不动——传统一线的引用标签也是半角，两边一致。
% 句点不用补：2015 样式本来就是半角加空格。圆括号与方括号同款，走
% \cs{bibleftparen}/\cs{bibrightparen} 领口。
%    \begin{macrocode}
\cs_new_protected:Npn \hit@bib@if@cjk@entry:TF #1#2
  {
    \ifboolexpr
      {
        test { \iffieldequalstr { userd } { chinese } } or
        test { \iffieldequalstr { userd } { japanese } } or
        test { \iffieldequalstr { userd } { korean } }
      }
      {#1} {#2}
  }
\cs_new_protected:Npn \hit@bib@patch@brackets:
  {
    \cs_if_exist:NF \__hit_bib_lbracket_std:
      {
        \cs_new_eq:NN \__hit_bib_lbracket_std: \bibleftbracket
        \cs_new_eq:NN \__hit_bib_rbracket_std: \bibrightbracket
      }
    \cs_set_protected:Npn \bibleftbracket
      {
        \str_if_eq:VnTF \g__hit_bib_punct_str { half }
          { \__hit_bib_lbracket_std: }
          {
            \str_if_eq:VnTF \g__hit_bib_punct_str { full }
              { \blx@postpunct ［ }
              {
                \hit@bib@if@cjk@entry:TF
                  { \blx@postpunct ［ } { \__hit_bib_lbracket_std: }
              }
          }
      }
    \cs_set_protected:Npn \bibrightbracket
      {
        \str_if_eq:VnTF \g__hit_bib_punct_str { half }
          { \__hit_bib_rbracket_std: }
          {
            \str_if_eq:VnTF \g__hit_bib_punct_str { full }
              { \blx@postpunct ］ \midsentence }
              {
                \hit@bib@if@cjk@entry:TF
                  { \blx@postpunct ］ \midsentence }
                  { \__hit_bib_rbracket_std: }
              }
          }
      }
    \cs_if_exist:NF \__hit_bib_lparen_std:
      {
        \cs_new_eq:NN \__hit_bib_lparen_std: \bibleftparen
        \cs_new_eq:NN \__hit_bib_rparen_std: \bibrightparen
      }
    \cs_set_protected:Npn \bibleftparen
      {
        \str_if_eq:VnTF \g__hit_bib_punct_str { half }
          { \__hit_bib_lparen_std: }
          {
            \str_if_eq:VnTF \g__hit_bib_punct_str { full }
              { \blx@postpunct （ }
              {
                \hit@bib@if@cjk@entry:TF
                  { \blx@postpunct （ } { \__hit_bib_lparen_std: }
              }
          }
      }
    \cs_set_protected:Npn \bibrightparen
      {
        \str_if_eq:VnTF \g__hit_bib_punct_str { half }
          { \__hit_bib_rparen_std: }
          {
            \str_if_eq:VnTF \g__hit_bib_punct_str { full }
              { \blx@postpunct ）\midsentence }
              {
                \hit@bib@if@cjk@entry:TF
                  { \blx@postpunct ）\midsentence }
                  { \__hit_bib_rparen_std: }
              }
          }
      }
  }
%    \end{macrocode}
% \end{macro}
%
% \begin{macro}{\hit@bib@patch@punct:}
% 逗号冒号分号的补丁，把 2025 代 \texttt{gbpunctwidth} 的能力补到 2015 代上。
% \pkg{biblatex} 的标点缓冲（重复合并、句点让位那套语义）最终从
% \cs{abx@puncthook}\marg{字符} 这一个口打出字形，补丁在这里按
% \option{punct} 换成全角，缓冲语义原封不动；句点问号叹号照旧半角。
% 全角标点后原样式还会发一个 \cs{addspace}（写在几十处调用点里），全角后
% 不该再空，打全角时立一个旗子，\cs{addspace}/\cs{addthinspace} 先照原样
% 冲刷缓冲（冲刷才会打出标点、立起旗子），再见旗吞空格、用后即清；旗子
% 每次进钩子先清一遍，不会窜到别的空格上。姓名之间的分隔符默认写的是裸
% \cs{space}（不走 \cs{addspace}，吞不到），把
% \texttt{multinamedelim}/\texttt{finalnamedelim} 重声明、
% \cs{pubdatadelim}/\cs{locnopubdelim} 重定义、期刊宏里那处内联的
% 卷号分隔用 \cs{patchcmd} 点补，全部换成走 \cs{addspace} 的写法，
% 半角档下排出来与原样一致。
%
% 标点缓冲在正文引用里也会走这条钩子（比如页码注的冒号），那边的全半角归
% \option{cite-punct} 管，所以钩子带 \cs{ifbibliography} 判据，只在文献表里
% 分发，引用侧走原样。
%
% 类型标与「不详」补白的方括号走的是 \cs{gbleftbracket} 一族——它同时也是
% 引用标签的括号，所以补丁带一道 \cs{ifbibliography} 判据：文献表里按
% \option{punct} 分发，引用侧按 \option{brackets-in-cite} 分发，
% 两个键各管各的。
%
% 挂载时机有讲究：\cs{abx@puncthook} 会被 \pkg{biblatex} 的语言初始化在
% \texttt{begindocument} 里重设（法文间距那套），补丁挂
% \texttt{begindocument/end}、注册在 \pkg{biblatex} 装载之后，排它后面。
%    \begin{macrocode}
\bool_new:N \g__hit_bib_gobble_bool
\cs_new_protected:Npn \hit@bib@punct@full:n #1
  {
    \str_case:nnF {#1}
      {
        { , } { ，\allowbreak \bool_gset_true:N \g__hit_bib_gobble_bool }
        { : } { ：\allowbreak \bool_gset_true:N \g__hit_bib_gobble_bool }
        { ; } { ；\allowbreak \bool_gset_true:N \g__hit_bib_gobble_bool }
      }
      { \__hit_bib_puncthook_std:n {#1} }
  }
\cs_new_protected:Npn \hit@bib@patch@punct:
  {
    \cs_if_exist:NF \__hit_bib_puncthook_std:n
      {
        \cs_new_eq:NN \__hit_bib_puncthook_std:n \abx@puncthook
        \cs_new_eq:NN \__hit_bib_addspace_std: \addspace
        \cs_new_eq:NN \__hit_bib_addthinspace_std: \addthinspace
      }
    \cs_set_protected:Npn \abx@puncthook ##1
      {
        \bool_gset_false:N \g__hit_bib_gobble_bool
        \ifbibliography
          {
            \str_if_eq:VnTF \g__hit_bib_punct_str { half }
              { \__hit_bib_puncthook_std:n {##1} }
              {
                \str_if_eq:VnTF \g__hit_bib_punct_str { full }
                  { \hit@bib@punct@full:n {##1} }
                  {
                    \hit@bib@if@cjk@entry:TF
                      { \hit@bib@punct@full:n {##1} }
                      { \__hit_bib_puncthook_std:n {##1} }
                  }
              }
          }
          { \__hit_bib_puncthook_std:n {##1} }
      }
    \cs_set_protected:Npn \addspace
      {
        \unspace \blx@postpunct
        \bool_if:NTF \g__hit_bib_gobble_bool
          { \bool_gset_false:N \g__hit_bib_gobble_bool }
          { \space }
        \blx@imc@resetpunctfont
      }
    \cs_set_protected:Npn \addthinspace
      {
        \unspace \blx@postpunct
        \bool_if:NTF \g__hit_bib_gobble_bool
          { \bool_gset_false:N \g__hit_bib_gobble_bool }
          { \hskip 0.16667 em \relax }
        \blx@imc@resetpunctfont
      }
    \DeclareDelimFormat { multinamedelim } { \addcomma \addspace }
    \DeclareDelimFormat { finalnamedelim } { \addcomma \addspace }
    \cs_if_exist:NT \compcitedelim
      {
        \cs_set_protected:Npn \compcitedelim  { \hit@bib@cite@delim: }
        \cs_set_protected:Npn \supercitedelim { \hit@bib@cite@delim: }
        \DeclareDelimFormat { multicitedelimiter } { \hit@bib@cite@delim: }
        \DeclareDelimFormat [ parencite ] { multicitedelimiter }
          { \hit@bib@cite@delim: }
      }
    \cs_if_exist:NT \pubdatadelim
      { \cs_set:Npn \pubdatadelim { \setunit* { \addcomma \addspace } } }
    \cs_if_exist:NT \locnopubdelim
      { \cs_set:Npn \locnopubdelim { \setunit* { \addcomma \addspace } } }
    \cs_if_exist:cT { abx@macro@journal+issuetitle }
      {
        \exp_args:Nc \patchcmd { abx@macro@journal+issuetitle }
          { \setunit* { \addcomma \space } }
          { \setunit* { \addcomma \addspace } }
          { } { }
      }
    \cs_if_exist:NF \__hit_bib_gblbracket_std:
      {
        \cs_new_eq:NN \__hit_bib_gblbracket_std: \gbleftbracket
        \cs_new_eq:NN \__hit_bib_gbrbracket_std: \gbrightbracket
      }
    \cs_set_protected:Npn \gbleftbracket
      {
        \ifbibliography
          { \hit@bib@punct@bracket:nN { ［ } \__hit_bib_gblbracket_std: }
          {
            \str_if_eq:VnTF \g__hit_bib_cite_brackets_str { full }
              { ［ } { \__hit_bib_gblbracket_std: }
          }
      }
    \cs_set_protected:Npn \gbrightbracket
      {
        \ifbibliography
          { \hit@bib@punct@bracket:nN { ］ } \__hit_bib_gbrbracket_std: }
          {
            \str_if_eq:VnTF \g__hit_bib_cite_brackets_str { full }
              { ］ } { \__hit_bib_gbrbracket_std: }
          }
      }
  }
\cs_new_protected:Npn \hit@bib@punct@bracket:nN #1#2
  {
    \str_if_eq:VnTF \g__hit_bib_punct_str { half }
      {#2}
      {
        \str_if_eq:VnTF \g__hit_bib_punct_str { full }
          {#1}
          { \hit@bib@if@cjk@entry:TF {#1} {#2} }
      }
  }
%    \end{macrocode}
% \end{macro}
%
% \begin{macro}{\hit@bib@cite@delim:,\hit@bib@declare@sourcemaps:}
% 引用序号分隔（biblatex 一线）：gb 系 cbx 把分隔收在 \cs{compcitedelim}、
% \cs{supercitedelim} 与 \texttt{multicitedelimiter} 三处，全指到这一个
% 分发宏，按 \option{cite-punct} 现值排，随设随生效。
%
% 数据清洗（biblatex 一线）：\file{.bst} 的 \texttt{convertpunct} 在排版期
% 对字段值做字符串替换，biber 这边同一层的东西是 sourcemap——装载时按
% \option{punct} 的档位声明一批替换步骤，替换对照表照抄 \file{.bst}。
% \texttt{by-lang} 档按条目判语言：题名含 CJK 判中文（正则与 gb 的 bbx 判
% \texttt{userd} 用的同一条），\texttt{final} 步骤不中就跳过整张 map。
% 只清洗文本字段，姓名字段不动（见键的说明）。步骤表先在配平的暂存
% token 表里攒好，再用 e 展开一次性组装进 \cs{DeclareSourcemap}。
% 装载时的档位定死，之后改 \option{punct} 只影响模板排的标点、
% 不重洗数据。
%    \begin{macrocode}
\cs_new_protected:Npn \hit@bib@cite@delim:
  {
    \str_case:VnF \g__hit_bib_cite_punct_str
      {
        { half-with-space } { \addcomma \addspace }
        { full } { \mbox{，} \allowbreak }
      }
      { \addcomma }
  }
\tl_new:N \l__hit_bib_steps_tl
\tl_new:N \l__hit_bib_maps_tl
\clist_const:Nn \c__hit_bib_map_fields_clist
  {
    title , subtitle , titleaddon , booktitle , journaltitle ,
    publisher , location , address , eventtitle , series
  }
\cs_new_protected:Npn \__hit_bib_full_steps:n #1
  {
    \tl_put_right:Nn \l__hit_bib_steps_tl
      {
        \step [ fieldsource = #1 , match = \regexp {．} , replace = \regexp {.~} ]
        \step [ fieldsource = #1 , match = \regexp {,\s+} , replace = \regexp {，} ]
        \step [ fieldsource = #1 , match = \regexp {:\s+} , replace = \regexp {：} ]
        \step [ fieldsource = #1 , match = \regexp {;\s+} , replace = \regexp {；} ]
        \step [ fieldsource = #1 , match = \regexp {!\s*} , replace = \regexp {！} ]
        \step [ fieldsource = #1 , match = \regexp {\?\s*} , replace = \regexp {？} ]
        \step [ fieldsource = #1 , match = \regexp {\s*\(} , replace = \regexp {（} ]
        \step [ fieldsource = #1 , match = \regexp {\)\s*} , replace = \regexp {）} ]
      }
  }
\cs_new_protected:Npn \__hit_bib_half_steps:n #1
  {
    \tl_put_right:Nn \l__hit_bib_steps_tl
      {
        \step [ fieldsource = #1 , match = \regexp {．\s*} , replace = \regexp {.~} ]
        \step [ fieldsource = #1 , match = \regexp {，\s*} , replace = \regexp {,~} ]
        \step [ fieldsource = #1 , match = \regexp {：\s*} , replace = \regexp {:~} ]
        \step [ fieldsource = #1 , match = \regexp {；\s*} , replace = \regexp {;~} ]
        \step [ fieldsource = #1 , match = \regexp {！\s*} , replace = \regexp {!~} ]
        \step [ fieldsource = #1 , match = \regexp {？\s*} , replace = \regexp {?~} ]
        \step [ fieldsource = #1 , match = \regexp {\s*（} , replace = \regexp {~(} ]
        \step [ fieldsource = #1 , match = \regexp {）\s*} , replace = \regexp {)~} ]
      }
  }
\cs_new_protected:Npn \__hit_bib_add_map:n #1
  {
    \tl_put_right:Nx \l__hit_bib_maps_tl
      { \exp_not:N \map [ overwrite ] { \exp_not:n {#1} \exp_not:V \l__hit_bib_steps_tl } }
  }
\cs_new_protected:Npn \hit@bib@declare@sourcemaps:
  {
    \tl_clear:N \l__hit_bib_maps_tl
    \str_case:VnF \g__hit_bib_punct_str
      {
        { half }
          {
            \tl_clear:N \l__hit_bib_steps_tl
            \clist_map_function:NN \c__hit_bib_map_fields_clist
              \__hit_bib_half_steps:n
            \__hit_bib_add_map:n { }
          }
        { by-lang }
          {
            \tl_clear:N \l__hit_bib_steps_tl
            \clist_map_function:NN \c__hit_bib_map_fields_clist
              \__hit_bib_full_steps:n
            \__hit_bib_add_map:n
              {
                \step [ fieldsource = title ,
                         match = \regexp { [\x{2FF0}-\x{9FA5}] } , final ]
              }
            \tl_clear:N \l__hit_bib_steps_tl
            \clist_map_function:NN \c__hit_bib_map_fields_clist
              \__hit_bib_half_steps:n
            \__hit_bib_add_map:n
              {
                \step [ fieldsource = title ,
                         notmatch = \regexp { [\x{2FF0}-\x{9FA5}] } , final ]
              }
          }
      }
      {
        \tl_clear:N \l__hit_bib_steps_tl
        \clist_map_function:NN \c__hit_bib_map_fields_clist
          \__hit_bib_full_steps:n
        \__hit_bib_add_map:n { }
      }
    \use:e
      {
        \exp_not:N \DeclareSourcemap
          {
            \exp_not:N \maps [ datatype = bibtex ]
              { \exp_not:V \l__hit_bib_maps_tl }
          }
      }
  }
%    \end{macrocode}
% \end{macro}
%
%    \begin{macrocode}
\cs_new_protected:Npn \__hit_bib_body_bibliography:n #1 { \printbibliography }
\cs_new_protected:Npn \__hit_bib_body_style:n #1
  { \msg_warning:nn { hithesis } { bib-style-ignored } }
\AddToHook { package / biblatex / before }
  { \cs_undefine:N \inlinecite }
\cs_new_protected:Npn \hit@bib@biblatex@common:
  {
    \cs_set:Npn \@inlinecite ##1 { \parencite {##1} }
    \cs_if_exist:NF \inlinecite
      { \DeclareRobustCommand \inlinecite { \@inlinecite } }
    \cs_set_eq:NN \onlinecite \inlinecite
    \cs_set:Npn \supercite { \cite }
    \cs_set_eq:NN \bibliographystyle \__hit_bib_body_style:n
    \AddToHook { begindocument / end }
      {
        \cs_set_eq:NN \bibliography \__hit_bib_body_bibliography:n
        \cs_set_eq:NN \bibliographystyle \__hit_bib_body_style:n
        \hit@bib@patch@punct:
      }
    \cs_if_exist:NT \mkgbnumlabel
      {
        \bool_if:NTF \g_bib_authorupper_bool
          { \ExecuteBibliographyOptions { gbnamefmt = uppercase } }
          { \ExecuteBibliographyOptions { gbnamefmt = lowercase } }
        \str_if_eq:VnT \g_bib_labelalign_str { left }
          { \ExecuteBibliographyOptions { gbalign = left } }
        \bool_if:NTF \g__hit_bib_slsn_bool
          { \ExecuteBibliographyOptions { gbpub = true } }
          { \ExecuteBibliographyOptions { gbpub = false } }
        \bool_if:NTF \g__hit_bib_showmark_bool
          { \ExecuteBibliographyOptions { gbtype = true } }
          { \ExecuteBibliographyOptions { gbtype = false } }
        \bool_if:NTF \g__hit_bib_showmedium_bool
          { \ExecuteBibliographyOptions { gbmedium = true } }
          { \ExecuteBibliographyOptions { gbmedium = false } }
        \bool_if:NTF \g__hit_bib_linktitle_bool
          { \ExecuteBibliographyOptions { gbtitlelink = true } }
          { \ExecuteBibliographyOptions { gbtitlelink = false } }
%    \end{macrocode}
%
% \option{sentence-case} 在这条线上改写 \pkg{biblatex-gb7714-2015} 留的三个钩子
% （\cs{mktitlecase} 一组，原样定义是恒等）。标志位在钩子里读，不是在这里读，
% 所以设得晚也照样生效，与 \file{.bst} 一线的时机对齐。
%
% 用带星号的 \cs{MakeSentenceCase*}：它只对判定为西文的条目动手，中文题名
% 不受影响。
%    \begin{macrocode}
        \cs_set:Npn \mktitlecase ##1
          {
            \bool_if:NTF \g__hit_bib_sentence_title_bool
              { \MakeSentenceCase* {##1} } {##1}
          }
        \cs_set:Npn \mkbooktitlecase ##1
          {
            \bool_if:NTF \g__hit_bib_sentence_booktitle_bool
              { \MakeSentenceCase* {##1} } {##1}
          }
        \cs_set:Npn \mkjournaltitlecase ##1
          {
            \bool_if:NTF \g__hit_bib_sentence_journal_bool
              { \MakeSentenceCase* {##1} } {##1}
          }
        \hit@bib@patch@brackets:
      }
    \bool_if:NTF \g__hit_bib_showurl_bool
      { \ExecuteBibliographyOptions { url = true } }
      { \ExecuteBibliographyOptions { url = false } }
    \bool_if:NTF \g__hit_bib_showdoi_bool
      { \ExecuteBibliographyOptions { doi = true } }
      { \ExecuteBibliographyOptions { doi = false } }
    \bool_if:NTF \g__hit_bib_showeprint_bool
      { \ExecuteBibliographyOptions { eprint = true } }
      { \ExecuteBibliographyOptions { eprint = false } }
    \bool_if:NT \g__hit_bib_correct_punct_bool
      { \hit@bib@declare@sourcemaps: }
    \int_compare:nNnTF \g_bib_maxauthor_int < 0
      { \ExecuteBibliographyOptions { maxbibnames = 999 , minbibnames = 999 } }
      {
        \use:e
          {
            \exp_not:N \ExecuteBibliographyOptions
              {
                maxbibnames = \int_use:N \g_bib_maxauthor_int ,
                minbibnames = \int_use:N \g_bib_maxauthor_int
              }
          }
      }
    \int_compare:nNnF \g__hit_bib_maxcite_int = { -2 }
      {
        \int_compare:nNnTF \g__hit_bib_maxcite_int < 0
          { \ExecuteBibliographyOptions { maxcitenames = 999 , mincitenames = 999 } }
          {
            \use:e
              {
                \exp_not:N \ExecuteBibliographyOptions
                  {
                    maxcitenames = \int_use:N \g__hit_bib_maxcite_int ,
                    mincitenames = \int_use:N \g__hit_bib_maxcite_int
                  }
              }
          }
      }
  }
\AddToHook { package / biblatex / after }
  {
    \hit@bib@biblatex@common:
    \cs_if_exist_use:N \hit@bib@biblatex@setup:
  }
%    \end{macrocode}
% \end{macro}
%
% \subsection{包冲突守卫}
%
% 「先装了 A 再装 B」才报错，所以一对包要两个方向各登记一条。用内核的
% \texttt{package/.../before} 与 \texttt{.../after} 钩子，不引 \pkg{filehook}。
% 类自己的装载不会撞上守卫：传统一线的装载在导言区末尾，且探过
% \pkg{biblatex} 不在才装。
%    \begin{macrocode}
\cs_new_protected:Npn \__hit_bib_conflict:nn #1#2
  {
    \AddToHook { package / #1 / after }
      {
        \AddToHook { package / #2 / before }
          { \msg_error:nnnn { hithesis } { package-conflict } {#1} {#2} }
      }
  }
\__hit_bib_conflict:nn { biblatex } { natbib }
\__hit_bib_conflict:nn { biblatex } { gbt7714 }
\__hit_bib_conflict:nn { biblatex } { cite }
\__hit_bib_conflict:nn { biblatex } { bibunits }
\__hit_bib_conflict:nn { biblatex } { chapterbib }
\__hit_bib_conflict:nn { biblatex } { multibib }
\__hit_bib_conflict:nn { natbib } { biblatex }
\__hit_bib_conflict:nn { gbt7714 } { biblatex }
\__hit_bib_conflict:nn { cite } { biblatex }
%    \end{macrocode}
%
% \subsection{文献库文件}
%
% \option{struct} 里的 \option{bibliography} 写的是给 \cs{bibliography} 的
% 文件名表（不带后缀）。\pkg{biblatex} 一线改喂 \cs{addbibresource}，一个
% 文件一句、后缀补上 \file{.bib}（写了后缀的不重复补）。挂导言区末尾做，
% \option{struct} 与 \option{backend} 谁先谁后都行。开关型取值
% （\texttt{false}/\texttt{true}/\texttt{auto}）不是文件名，跳过。
%    \begin{macrocode}
\cs_new_protected:Npn \__hit_bib_add_resource:n #1
  {
    \str_if_eq:eeTF { \str_range:nnn {#1} {-4} {-1} } { .bib }
      { \addbibresource {#1} }
      { \addbibresource { #1.bib } }
  }
\cs_new_protected:Npn \__hit_bib_resources:n #1
  {
    \str_case:nnF {#1}
      {
        { }      { \msg_warning:nn { hithesis } { bib-no-resource } }
        { false } { }
        { true }  { \msg_warning:nn { hithesis } { bib-no-resource } }
        { auto }  { \msg_warning:nn { hithesis } { bib-no-resource } }
      }
      { \clist_map_inline:nn {#1} { \__hit_bib_add_resource:n {##1} } }
  }
\AddToHook { begindocument / before }
  {
    \@ifpackageloaded { biblatex }
      { \exp_args:Nv \__hit_bib_resources:n { g__hit_structure_bibliography_tl } }
      { }
  }
%</shared-bibstyle>
%    \end{macrocode}
%
% \Finale
